\section{Experiments and Observations}
\label{sec:experiments}

All numbers below come from a single recorded run of the frozen configuration at commit
\texttt{\res{commit}}. That run covers $3$ conditions $\times$ \res{nseeds} world seeds $\times$
\res{norders} orders $\times$ $8$ arms, and \res{ncellsok} of its \res{ncells} cells completed. The
cells are not independent replicates: orders are nested in world seeds, and the three conditions reuse the
same worlds. The whole run used \res{cpu}\,s of a \res{cpubudget}\,s CPU budget, with peak resident
memory \res{rss}\,MiB, on one worker and one thread.

\paragraph{Decision.} The pre-registered rule returned \emph{order penalty added utility not supported},
with \res{nreasons} failed conditions. The candidate's stationary gain over no adaptation was
\res{ev.gain.r0}, \res{ev.gain.r1} and \res{ev.gain.r2}\,pp in the three seeds, so the gain condition
failed. Its SD reduction relative to EM-toy was \res{ev.stdred.r0}, \res{ev.stdred.r1} and
\res{ev.stdred.r2}\,pp, negative in one seed, so the robustness condition failed. It dominated none of the
five required controls. In drift its seed-mean gain over no adaptation was \res{ev.drift.re.gain}\,pp
(current regime) and \res{ev.drift.on.gain}\,pp (online).

% AUTO-GENERATED by paper/scripts/export_results.py. Do not edit.
\begin{table}[t]
\centering
\small
\setlength{\tabcolsep}{3.2pt}
\caption{Primary stationary condition (fixed batch membership, batch order permuted). Err: mean terminal common-holdout error over the 8 observed orders (\%); SD: standard deviation over those orders (percentage points, pp); Gain: no-adaptation error minus arm error (pp; positive is better). Orders are nested in world seeds and are not independent replicates.}
\label{tab:stationary}
\begin{tabular}{lrrr@{\hspace{7pt}}rrr@{\hspace{7pt}}rrr@{\hspace{7pt}}}
\toprule
 & \multicolumn{3}{c}{World seed r0} & \multicolumn{3}{c}{World seed r1} & \multicolumn{3}{c}{World seed r2} \\
\cmidrule(lr){2-4} \cmidrule(lr){5-7} \cmidrule(lr){8-10}
Arm & Err & SD & Gain & Err & SD & Gain & Err & SD & Gain \\
\midrule
No adaptation & \res{stat.te.noad.mean.r0} & \res{stat.te.noad.std.r0} & \res{stat.te.noad.gain.r0} & \res{stat.te.noad.mean.r1} & \res{stat.te.noad.std.r1} & \res{stat.te.noad.gain.r1} & \res{stat.te.noad.mean.r2} & \res{stat.te.noad.std.r2} & \res{stat.te.noad.gain.r2} \\
EM-toy (full affine) & \res{stat.te.full.mean.r0} & \res{stat.te.full.std.r0} & \res{stat.te.full.gain.r0} & \res{stat.te.full.mean.r1} & \res{stat.te.full.std.r1} & \res{stat.te.full.gain.r1} & \res{stat.te.full.mean.r2} & \res{stat.te.full.std.r2} & \res{stat.te.full.gain.r2} \\
EM-toy, lr $\times$0.3 & \res{stat.te.lrA.mean.r0} & \res{stat.te.lrA.std.r0} & \res{stat.te.lrA.gain.r0} & \res{stat.te.lrA.mean.r1} & \res{stat.te.lrA.std.r1} & \res{stat.te.lrA.gain.r1} & \res{stat.te.lrA.mean.r2} & \res{stat.te.lrA.std.r2} & \res{stat.te.lrA.gain.r2} \\
EM-toy, lr $\times$0.1 & \res{stat.te.lrB.mean.r0} & \res{stat.te.lrB.std.r0} & \res{stat.te.lrB.gain.r0} & \res{stat.te.lrB.mean.r1} & \res{stat.te.lrB.std.r1} & \res{stat.te.lrB.gain.r1} & \res{stat.te.lrB.mean.r2} & \res{stat.te.lrB.std.r2} & \res{stat.te.lrB.gain.r2} \\
Order-penalized subspace (candidate) & \res{stat.te.cand.mean.r0} & \res{stat.te.cand.std.r0} & \res{stat.te.cand.gain.r0} & \res{stat.te.cand.mean.r1} & \res{stat.te.cand.std.r1} & \res{stat.te.cand.gain.r1} & \res{stat.te.cand.mean.r2} & \res{stat.te.cand.std.r2} & \res{stat.te.cand.gain.r2} \\
Utility-only subspace & \res{stat.te.util.mean.r0} & \res{stat.te.util.std.r0} & \res{stat.te.util.gain.r0} & \res{stat.te.util.mean.r1} & \res{stat.te.util.std.r1} & \res{stat.te.util.gain.r1} & \res{stat.te.util.mean.r2} & \res{stat.te.util.std.r2} & \res{stat.te.util.gain.r2} \\
Per-step norm-matched (causal) & \res{stat.te.nstep.mean.r0} & \res{stat.te.nstep.std.r0} & \res{stat.te.nstep.gain.r0} & \res{stat.te.nstep.mean.r1} & \res{stat.te.nstep.std.r1} & \res{stat.te.nstep.gain.r1} & \res{stat.te.nstep.mean.r2} & \res{stat.te.nstep.std.r2} & \res{stat.te.nstep.gain.r2} \\
Global norm-matched (calibrated) & \res{stat.te.nglob.mean.r0} & \res{stat.te.nglob.std.r0} & \res{stat.te.nglob.gain.r0} & \res{stat.te.nglob.mean.r1} & \res{stat.te.nglob.std.r1} & \res{stat.te.nglob.gain.r1} & \res{stat.te.nglob.mean.r2} & \res{stat.te.nglob.std.r2} & \res{stat.te.nglob.gain.r2} \\
\bottomrule
\end{tabular}
\end{table}


\paragraph{Primary stationary condition (Table~\ref{tab:stationary}).} With fixed batch membership, EM-toy did
not improve on no adaptation in any seed: its gains were \res{stat.te.full.gain.r0},
\res{stat.te.full.gain.r1} and \res{stat.te.full.gain.r2}\,pp. Its order SD was
\res{stat.te.full.std.r0}, \res{stat.te.full.std.r1} and \res{stat.te.full.std.r2}\,pp. The small-step
controls had a smaller SD than EM-toy in every seed. For lr$\times$0.1 the gains were
\res{stat.te.lrB.gain.r0}, \res{stat.te.lrB.gain.r1} and \res{stat.te.lrB.gain.r2}\,pp, all below the
MIE. The candidate had the lowest error in seed r0, tied with the identical utility-only run discussed below,
but it was worse than no adaptation in seeds r1 and r2. In seed r2 its
error was \res{stat.te.cand.mean.r2}\% against \res{stat.te.noad.mean.r2}\% without adaptation. In this
environment low order sensitivity is attained most easily by adapting little: no adaptation has zero SD by
construction.

% AUTO-GENERATED by paper/scripts/export_results.py. Do not edit.
\begin{table}[t]
\centering
\small
\setlength{\tabcolsep}{3.2pt}
\caption{Paired differences, candidate minus control, computed per world seed over the same 8 orders (pp; negative favours the candidate). $\Delta$Err: mean terminal error difference; $\Delta$SD: difference of across-order standard deviations. Primary stationary condition.}
\label{tab:paired}
\begin{tabular}{lrr@{\hspace{7pt}}rr@{\hspace{7pt}}rr@{\hspace{7pt}}}
\toprule
 & \multicolumn{2}{c}{World seed r0} & \multicolumn{2}{c}{World seed r1} & \multicolumn{2}{c}{World seed r2} \\
\cmidrule(lr){2-3} \cmidrule(lr){4-5} \cmidrule(lr){6-7}
Control & $\Delta$Err & $\Delta$SD & $\Delta$Err & $\Delta$SD & $\Delta$Err & $\Delta$SD \\
\midrule
Utility-only subspace & \res{stat.te.pair.util.dmean.r0} & \res{stat.te.pair.util.dstd.r0} & \res{stat.te.pair.util.dmean.r1} & \res{stat.te.pair.util.dstd.r1} & \res{stat.te.pair.util.dmean.r2} & \res{stat.te.pair.util.dstd.r2} \\
EM-toy, lr $\times$0.3 & \res{stat.te.pair.lrA.dmean.r0} & \res{stat.te.pair.lrA.dstd.r0} & \res{stat.te.pair.lrA.dmean.r1} & \res{stat.te.pair.lrA.dstd.r1} & \res{stat.te.pair.lrA.dmean.r2} & \res{stat.te.pair.lrA.dstd.r2} \\
EM-toy, lr $\times$0.1 & \res{stat.te.pair.lrB.dmean.r0} & \res{stat.te.pair.lrB.dstd.r0} & \res{stat.te.pair.lrB.dmean.r1} & \res{stat.te.pair.lrB.dstd.r1} & \res{stat.te.pair.lrB.dmean.r2} & \res{stat.te.pair.lrB.dstd.r2} \\
Per-step norm-matched (causal) & \res{stat.te.pair.nstep.dmean.r0} & \res{stat.te.pair.nstep.dstd.r0} & \res{stat.te.pair.nstep.dmean.r1} & \res{stat.te.pair.nstep.dstd.r1} & \res{stat.te.pair.nstep.dmean.r2} & \res{stat.te.pair.nstep.dstd.r2} \\
Global norm-matched (calibrated) & \res{stat.te.pair.nglob.dmean.r0} & \res{stat.te.pair.nglob.dstd.r0} & \res{stat.te.pair.nglob.dmean.r1} & \res{stat.te.pair.nglob.dstd.r1} & \res{stat.te.pair.nglob.dmean.r2} & \res{stat.te.pair.nglob.dstd.r2} \\
EM-toy (full affine) & \res{stat.te.pair.full.dmean.r0} & \res{stat.te.pair.full.dstd.r0} & \res{stat.te.pair.full.dmean.r1} & \res{stat.te.pair.full.dstd.r1} & \res{stat.te.pair.full.dmean.r2} & \res{stat.te.pair.full.dstd.r2} \\
\bottomrule
\end{tabular}
\end{table}


\paragraph{Paired comparison with the controls (Table~\ref{tab:paired}).} Relative to the utility-only
control, the candidate's mean-error difference was \res{stat.te.pair.util.dmean.r0},
\res{stat.te.pair.util.dmean.r1} and \res{stat.te.pair.util.dmean.r2}\,pp. Relative to the
reduced-learning-rate and norm-matched controls it was better in seed r0 and worse in seeds r1 and r2. Under
the rule, a control counts as explained away only if the candidate beats it with the same sign in all three
seeds; no control met that criterion.

% AUTO-GENERATED by paper/scripts/export_results.py. Do not edit.
\begin{table}[t]
\centering
\small
\setlength{\tabcolsep}{3pt}
\caption{Candidate (order-penalized) versus utility-only subspace per condition and world seed. Pool indices of the selected directions (both fitters select from the same shared pool and meta statistics), tuned learning rates, and whether the terminal parameters (SHA-256 of $\theta_T$) and every per-step update norm coincide on all 8 orders, read from the raw logs.}
\label{tab:identity}
\begin{tabular}{llllrrcc}
\toprule
Condition & Seed & Cand.\ idx & Util.\ idx & lr cand. & lr util. & same $\theta_T$ & same steps \\
\midrule
stationary & r0 & [4, 3] & [4, 3] & 0.1 & 0.1 & yes & yes \\
stationary & r1 & [1, 5] & [0, 2] & 2 & 0.05 & no & no \\
stationary & r2 & [3, 5] & [3, 5] & 2 & 2 & yes & yes \\
drift & r0 & [4, 3] & [4, 3] & 2 & 2 & yes & yes \\
drift & r1 & [1, 5] & [0, 2] & 2 & 0.2 & no & no \\
drift & r2 & [3, 5] & [3, 5] & 2 & 2 & yes & yes \\
reshuffle & r0 & [4, 3] & [4, 3] & 0.1 & 0.1 & yes & yes \\
reshuffle & r1 & [1, 5] & [0, 2] & 2 & 0.05 & no & no \\
reshuffle & r2 & [3, 5] & [3, 5] & 2 & 2 & yes & yes \\
\bottomrule
\end{tabular}
\end{table}


\paragraph{Why the candidate and the utility-only control coincide in two seeds (Table~\ref{tab:identity}).}
The candidate's greedy early stop selected rank \res{ranks} in every condition and seed. The utility-only
control is rank-matched and reuses the same pool and meta statistics, so it selects the top two directions by
supervised gain. In seeds r0 and r2 these were exactly the directions the penalized score kept, so the two
arms used the same projector $P_SP_S^\top$. Development tuning then selected the same learning rate for
both, because tuning sees identical computations. Everything else is shared by construction: the source
model and $\theta_0$ (checked before every run), the batches and their order, and a deterministic update
without randomness, momentum or buffers. The two runs are therefore the same computation. The raw logs confirm
identical terminal parameters and identical per-step update norms on all orders in \res{nidentical} of the
\res{nidentityrows} condition--seed pairs. The order penalty altered the selected subspace only in seed r1,
where the candidate kept pool directions $[1,5]$ instead of $[0,2]$, development tuning selected a larger
learning rate, and the candidate was worse by \res{stat.te.pair.util.dmean.r1}\,pp. We do not read efficacy
into any single seed. The observation is that, in this implementation, the penalty rarely changed the
selection, and when it did the change did not help.

% AUTO-GENERATED by paper/scripts/export_results.py. Do not edit.
\begin{table}[t]
\centering
\small
\setlength{\tabcolsep}{3.2pt}
\caption{Drift condition (single-domain batches fixed; orders permute the domain-block sequence and batches within blocks). Err: mean over the four block ends of the current-regime holdout error (\%); SD over orders (pp); Gain vs.\ no adaptation (pp).}
\label{tab:drift}
\begin{tabular}{lrrr@{\hspace{7pt}}rrr@{\hspace{7pt}}rrr@{\hspace{7pt}}}
\toprule
 & \multicolumn{3}{c}{World seed r0} & \multicolumn{3}{c}{World seed r1} & \multicolumn{3}{c}{World seed r2} \\
\cmidrule(lr){2-4} \cmidrule(lr){5-7} \cmidrule(lr){8-10}
Arm & Err & SD & Gain & Err & SD & Gain & Err & SD & Gain \\
\midrule
No adaptation & \res{drift.re.noad.mean.r0} & \res{drift.re.noad.std.r0} & \res{drift.re.noad.gain.r0} & \res{drift.re.noad.mean.r1} & \res{drift.re.noad.std.r1} & \res{drift.re.noad.gain.r1} & \res{drift.re.noad.mean.r2} & \res{drift.re.noad.std.r2} & \res{drift.re.noad.gain.r2} \\
EM-toy (full affine) & \res{drift.re.full.mean.r0} & \res{drift.re.full.std.r0} & \res{drift.re.full.gain.r0} & \res{drift.re.full.mean.r1} & \res{drift.re.full.std.r1} & \res{drift.re.full.gain.r1} & \res{drift.re.full.mean.r2} & \res{drift.re.full.std.r2} & \res{drift.re.full.gain.r2} \\
EM-toy, lr $\times$0.3 & \res{drift.re.lrA.mean.r0} & \res{drift.re.lrA.std.r0} & \res{drift.re.lrA.gain.r0} & \res{drift.re.lrA.mean.r1} & \res{drift.re.lrA.std.r1} & \res{drift.re.lrA.gain.r1} & \res{drift.re.lrA.mean.r2} & \res{drift.re.lrA.std.r2} & \res{drift.re.lrA.gain.r2} \\
EM-toy, lr $\times$0.1 & \res{drift.re.lrB.mean.r0} & \res{drift.re.lrB.std.r0} & \res{drift.re.lrB.gain.r0} & \res{drift.re.lrB.mean.r1} & \res{drift.re.lrB.std.r1} & \res{drift.re.lrB.gain.r1} & \res{drift.re.lrB.mean.r2} & \res{drift.re.lrB.std.r2} & \res{drift.re.lrB.gain.r2} \\
Order-penalized subspace (candidate) & \res{drift.re.cand.mean.r0} & \res{drift.re.cand.std.r0} & \res{drift.re.cand.gain.r0} & \res{drift.re.cand.mean.r1} & \res{drift.re.cand.std.r1} & \res{drift.re.cand.gain.r1} & \res{drift.re.cand.mean.r2} & \res{drift.re.cand.std.r2} & \res{drift.re.cand.gain.r2} \\
Utility-only subspace & \res{drift.re.util.mean.r0} & \res{drift.re.util.std.r0} & \res{drift.re.util.gain.r0} & \res{drift.re.util.mean.r1} & \res{drift.re.util.std.r1} & \res{drift.re.util.gain.r1} & \res{drift.re.util.mean.r2} & \res{drift.re.util.std.r2} & \res{drift.re.util.gain.r2} \\
Per-step norm-matched (causal) & \res{drift.re.nstep.mean.r0} & \res{drift.re.nstep.std.r0} & \res{drift.re.nstep.gain.r0} & \res{drift.re.nstep.mean.r1} & \res{drift.re.nstep.std.r1} & \res{drift.re.nstep.gain.r1} & \res{drift.re.nstep.mean.r2} & \res{drift.re.nstep.std.r2} & \res{drift.re.nstep.gain.r2} \\
Global norm-matched (calibrated) & \res{drift.re.nglob.mean.r0} & \res{drift.re.nglob.std.r0} & \res{drift.re.nglob.gain.r0} & \res{drift.re.nglob.mean.r1} & \res{drift.re.nglob.std.r1} & \res{drift.re.nglob.gain.r1} & \res{drift.re.nglob.mean.r2} & \res{drift.re.nglob.std.r2} & \res{drift.re.nglob.gain.r2} \\
\bottomrule
\end{tabular}
\end{table}


\paragraph{Drift (Table~\ref{tab:drift}).} Scored on the current regime's holdout at each block end, no
adapting arm improved on no adaptation by the MIE in any seed. The candidate's seed-mean difference from the
utility-only control was \res{ev.drift.re.vsutil}\,pp for the current regime (within the tolerance of
\res{tolg}\,pp) and \res{ev.drift.on.vsutil}\,pp for online error (outside it). Its gain over no
adaptation was negative on both metrics. Online errors and the mixture-holdout errors are in
Tables~\ref{tab:drift-online} and~\ref{tab:drift-mixture}. Drift results are reported apart from stationary
results and are not pooled with them.

\paragraph{Secondary condition and diagnostics.} Re-forming batches for every order
(Table~\ref{tab:reshuffle}) gave the same qualitative pattern: the candidate-minus-utility-only difference
was \res{reshuf.te.pair.util.dmean.r0}, \res{reshuf.te.pair.util.dmean.r1} and
\res{reshuf.te.pair.util.dmean.r2}\,pp.

The candidate's tuned learning rate was the largest grid value, \res{lrmax}, in \res{ncandbound} of $9$
condition--seed pairs. We therefore cannot claim that its learning rate was well tuned, and we did not widen
the grid.

Predicted classes on the holdout became more concentrated in seed r2. The largest class share was
\res{stat.diag.cand.maxshare.r2} for the candidate against \res{stat.diag.noad.maxshare.r2} without
adaptation, and \res{drift.diag.cand.maxshare.r2} in drift. This is partial concentration, not a
single-class collapse (Table~\ref{tab:diag}).

Actual per-arm spend, including the lockstep reference of the causal control and the shared meta statistics
of the utility-only control, is listed in Table~\ref{tab:cost}.
